\section{李代数理论}

一旦变换群与李代数之间的对应建立起来，这一理论就将明显地转向更加代数的方向，并以对李代数的深入研究为中心。\(^{13}\)

从 1888 到 1894 年的第一个短暂时期，以恩格尔、他的学生乌姆劳夫以及特别是基灵与 E. 嘉当的工作为标志，以关于复李代数的一系列引人注目的结果而告终。我们在上面已经看到，可解李代数的概念要归功于李本人，他（在复的情形）证明了可解线性李代数化为三角形约化定理（{[}203{]}，v. I, p.~270）。\(^{14}\) 基灵观察到{[}180{]}，李代数中存在一个最大的可解理想（今天称为根），且李代数对其根的商的根为零；他把根为零的李代数称为半单的，并证明它们是单代数的积（后一概念已经由

李引入，李还证明了“经典”李代数的单性（{[}203{]}，v. 3, p.~682））。

另一方面，基灵在李代数中引入了特征方程 \(\det(\operatorname{ad}(x)-\omega.1)=0\)，李在研究包含李代数中一个给定元素的二维李子代数时已经遇到过它。关于这些方法的分析——基灵以透彻的方式研究半单代数的“一般的”特征方程的根的性质，最终得到他最值得注意的结果，即单（复）李代数的完全确定——我们请读者参看其他历史注记。\(^{15}\)

基灵证明，可解代数的导代数“秩为 0”（这意味着 ad \(x\) 对该代数的每个元素 \(x\) 都是幂零的）。不久之后，恩格尔证明了“秩为 0”的代数是可解的（这一论断实质上就是今天所谓的恩格尔定理）。另一方面，E. 嘉当在他的学位论文中引入了今天所谓的“基灵型”，并建立了借助这个型刻画可解李代数与半单李代数的两个基本判据。

基灵曾断言（{[}180{]}，IV），李代数的导代数是一个半单代数与其幂零的根之和，但他的证明是不完全的。稍后，E. 嘉当未加证明地宣布（{[}52 a{]}，v. I₁, p.~104），更一般地，每个李代数都是其根与一个半单子代数之和；这个方向上当时无可争议地确立的唯一结果，是恩格尔的一个定理，它断言在每个非可解李代数中存在一个三维单李子代数。嘉当论断的第一个公开发表的证明（对复李代数）属于 E. E. 列维{[}200{]}；另一个证明（对实情形同样有效）由 J. H. C. 怀特海于 1936 年给出{[}334 a{]}。1942 年，A. 马尔采夫以关于“列维截口”在共轭意义下的唯一性定理完善了这一结果。

早在最初的工作中，李就给自己提出了每个李代数与一个线性李代数同构的问题。他曾以为通过考虑伴随表示就肯定地解决了它（并由此推出他的“第三定理”的一个证明）（{[}202{]}，vol.~5, Abh. III, pp.~42-75）；他很快认识到他的证明只对中心为零的李代数正确；在他之后，这个问题长时间悬而未决，直到 1935 年才由阿多肯定地解决{[}2 a{]}。另一方面，李实质上还给自己提出过确定单李代数的最小维数线性表示的问题，并对经典代数解决了它；在他的学位论文中，嘉当还对例外单代数解决了这个问题；\(^{16}\) 他在此背景下使用的方法，二十年后将被他推广，以得到实的或复的单李代数的所有不可约表示。

线性表示的完全可约性这一性质，似乎最早是由施图迪（以几何的形式）遇到的。在一份未发表但被（{[}203{]}，v. 3, pp.~785-788）引用的手稿中，他对 SL(2, C) 的李代数的线性表示证明了这一性质，并对 SL(3, C) 与 SL(4, C) 得到了部分结果。李与恩格尔借此机会猜想，完全可约性定理对无论什么 n 的 SL(n, C) 都成立。半单李代数的线性表示的完全可约性由 H. 外尔于 1925 年\(^{17}\)通过一个整体性论证（见后文）建立。第一个代数证明由卡西米尔与范德瓦尔登于 1935 年得到{[}55{]}；其后 R. 布饶尔{[}34 b{]}与 J. H. C. 怀特海{[}334 b{]}给出了另一些代数证明。

最后，在对指数映射的研究（参见后文）中，H. 庞加莱（{[}251{]}，v. III）考虑由一个李代数的算子生成的所有阶微分算子的结合代数；他实质上证明：若 \((X_{i})_{1\leq i\leq n}\) 是李代数的一组基，则由诸 \(X_{i}\) 生成的结合代数以诸 \(X_{i}\) 的某些对称函数（由一个给定单项式通过其因子的所有置换得到的非交换“单项式”之和）为基。他的证明的本质部分是代数性的，使他能够得到包络代数的结构。类似的证明由 G. 伯克霍夫{[}22 b{]}与 E. 维特{[}337 b{]}于 1937 年给出。\(^{18}\)

上面引用的工作大多限于实的或复的李代数，只有它们才对应于通常意义下的李群。对 R 或 C 以外的域上的李代数的研究由雅各布森{[}172 a{]}着手进行，他证明了经典结果的大部分在特征零的域上仍然成立。

